\documentclass{meetingminutes}
\meetingcommittee{คณะกรรมการบัณฑิตศึกษาประจำมหาวิทยาลัยราชภัฏอุตรดิตถ์}
\meetingnumber{๒}{๒๕๖๗}
\meetingdate{วันที่ ๔ มีนาคม พ.ศ. ๒๕๖๗}
\meetingplace{ห้องประชุมมหาวิทยาลัยราชภัฏอุตรดิตถ์}
\meetingstart{๑๓.๐๐ น.}
\meetingend{๑๖.๐๐ น.}

\begin{document}
\maketitle

\psection{ผู้มาประชุม}
\begin{participants}
  \pmember{1}{ประธานคณะกรรมการบัณฑิตศึกษา}{มหาวิทยาลัยราชภัฏอุตรดิตถ์}{ประธาน}
  \pmember{2}{ผู้แทนคณะเทคโนโลยีอุตสาหกรรม}{คณะเทคโนโลยีอุตสาหกรรม}{กรรมการ}
  \pmember{3}{ผู้แทนหลักสูตร วศ.ม. วิศวกรรมคอมพิวเตอร์และปัญญาประดิษฐ์}{คณะเทคโนโลยีอุตสาหกรรม}{กรรมการ}
  \pmember{4}{กรรมการบัณฑิตศึกษาประจำมหาวิทยาลัย}{มหาวิทยาลัยราชภัฏอุตรดิตถ์}{กรรมการ}
\end{participants}

\psection{ผู้ไม่มาประชุม}
\noindent ไม่มี\par

\startmeeting
\openingchair{ประธานคณะกรรมการบัณฑิตศึกษา}{กล่าวเปิดการประชุมและดำเนินการประชุมตามระเบียบวาระ ดังนี้}

% ============================================================
\agenda{๑}{เรื่องแจ้งเพื่อทราบ}
% ============================================================

ประธานแจ้งให้ที่ประชุมทราบว่า มีหลักสูตรระดับบัณฑิตศึกษาเสนอขอรับความเห็นชอบในการประชุมครั้งนี้
รวมทั้งหลักสูตรวิศวกรรมศาสตรมหาบัณฑิต สาขาวิชาวิศวกรรมคอมพิวเตอร์และปัญญาประดิษฐ์
คณะเทคโนโลยีอุตสาหกรรม

\noindent\textbf{มติที่ประชุม}\quad ที่ประชุมรับทราบ

% ============================================================
\agenda{๒}{พิจารณาให้ความเห็นชอบหลักสูตรวิศวกรรมศาสตรมหาบัณฑิต
สาขาวิชาวิศวกรรมคอมพิวเตอร์และปัญญาประดิษฐ์}
% ============================================================

ผู้แทนหลักสูตรนำเสนอ มคอ.๒ ฉบับสมบูรณ์ต่อที่ประชุมคณะกรรมการบัณฑิตศึกษา
ครอบคลุมสาระสำคัญดังนี้

\begin{thailist}
  \item ชื่อหลักสูตรและปริญญา: วิศวกรรมศาสตรมหาบัณฑิต สาขาวิชาวิศวกรรมคอมพิวเตอร์และปัญญาประดิษฐ์
  \item ผลลัพธ์การเรียนรู้ที่คาดหวัง (PLO) จำนวน ๕ ข้อ
  \item โครงสร้างหลักสูตร ๒ แผน ได้แก่ แผน ก (วิทยานิพนธ์) และแผน ข (สารนิพนธ์)
  \item คุณสมบัติอาจารย์ประจำหลักสูตร ๗ คน ตรงตามเกณฑ์ กมอ. พ.ศ. ๒๕๖๕
  (ระดับผู้ช่วยศาสตราจารย์/รองศาสตราจารย์ และสำเร็จการศึกษาระดับปริญญาเอก)
  \item เกณฑ์รับนักศึกษาและแผนการเปิดรับ
\end{thailist}

ที่ประชุมร่วมกันพิจารณาและอภิปราย ไม่มีข้อสังเกตเพิ่มเติม

\resolution{ที่ประชุมมีมติให้ความเห็นชอบ มคอ.๒ หลักสูตรวิศวกรรมศาสตรมหาบัณฑิต
สาขาวิชาวิศวกรรมคอมพิวเตอร์และปัญญาประดิษฐ์ และให้เสนอต่อสภาวิชาการเพื่อพิจารณาต่อไป}

% ============================================================
\agenda{๓}{เรื่องอื่น ๆ}
% ============================================================

ไม่มีวาระเพิ่มเติม

\noindent\textbf{มติที่ประชุม}\quad ที่ประชุมรับทราบ

\adjournmeeting

\begin{sigblock}
\typist{ผู้แทนหลักสูตร วศ.ม. วิศวกรรมคอมพิวเตอร์และปัญญาประดิษฐ์}
\reviewer{ผู้แทนคณะเทคโนโลยีอุตสาหกรรม}{กรรมการ}
\chairsig{ประธานคณะกรรมการบัณฑิตศึกษา}{ประธานที่ประชุม}
\end{sigblock}

\end{document}
